\documentclass[10pt,a4paper]{amsart}
\usepackage{amsmath,amssymb,amsthm,graphicx}
% Retained minimal preamble for the Stacks Project chapter varieties.tex.
% Verbatim excerpt of the chapter preamble: Xy-pic setup (native preamble lines 14-18), the native theorem
% declarations with the mathematical macros (native preamble lines 154-217), and the project own no-op
% definitions of the reference/slogan/history display environments (native wrapper.tex lines 1-3) that the
% counted statement uses for its slogan. Copied mechanically; nothing added, removed or redefined.
% For commutative diagrams we use Xy-pic
\usepackage[all]{xy}

% We use 2cell for 2-commutative diagrams.
\xyoption{2cell}
% Theorem environments.
%
\theoremstyle{plain}
\newtheorem{theorem}[subsection]{Theorem}
\newtheorem{proposition}[subsection]{Proposition}
\newtheorem{lemma}[subsection]{Lemma}

\theoremstyle{definition}
\newtheorem{definition}[subsection]{Definition}
\newtheorem{example}[subsection]{Example}
\newtheorem{exercise}[subsection]{Exercise}
\newtheorem{situation}[subsection]{Situation}

\theoremstyle{remark}
\newtheorem{remark}[subsection]{Remark}
\newtheorem{remarks}[subsection]{Remarks}

\numberwithin{equation}{subsection}

% Macros
%
\def\lim{\mathop{\mathrm{lim}}\nolimits}
\def\colim{\mathop{\mathrm{colim}}\nolimits}
\def\Spec{\mathop{\mathrm{Spec}}}
\def\Hom{\mathop{\mathrm{Hom}}\nolimits}
\def\Ext{\mathop{\mathrm{Ext}}\nolimits}
\def\SheafHom{\mathop{\mathcal{H}\!\mathit{om}}\nolimits}
\def\SheafExt{\mathop{\mathcal{E}\!\mathit{xt}}\nolimits}
\def\Sch{\mathit{Sch}}
\def\Mor{\mathop{\mathrm{Mor}}\nolimits}
\def\Ob{\mathop{\mathrm{Ob}}\nolimits}
\def\Sh{\mathop{\mathit{Sh}}\nolimits}
\def\NL{\mathop{N\!L}\nolimits}
\def\CH{\mathop{\mathrm{CH}}\nolimits}
\def\proetale{{pro\text{-}\acute{e}tale}}
\def\etale{{\acute{e}tale}}
\def\QCoh{\mathit{QCoh}}
\def\Ker{\mathop{\mathrm{Ker}}}
\def\Im{\mathop{\mathrm{Im}}}
\def\Coker{\mathop{\mathrm{Coker}}}
\def\Coim{\mathop{\mathrm{Coim}}}

% Boxtimes
%
\DeclareMathSymbol{\boxtimes}{\mathbin}{AMSa}{"02}

%
% Macros for moduli stacks/spaces
%
\def\QCohstack{\mathcal{QC}\!\mathit{oh}}
\def\Cohstack{\mathcal{C}\!\mathit{oh}}
\def\Spacesstack{\mathcal{S}\!\mathit{paces}}
\def\Quotfunctor{\mathrm{Quot}}
\def\Hilbfunctor{\mathrm{Hilb}}
\def\Curvesstack{\mathcal{C}\!\mathit{urves}}
\def\Polarizedstack{\mathcal{P}\!\mathit{olarized}}
\def\Complexesstack{\mathcal{C}\!\mathit{omplexes}}
% \Pic is the operator that assigns to X its picard group, usage \Pic(X)
% \Picardstack_{X/B} denotes the Picard stack of X over B
% \Picardfunctor_{X/B} denotes the Picard functor of X over B
\def\Pic{\mathop{\mathrm{Pic}}\nolimits}
\def\Picardstack{\mathcal{P}\!\mathit{ic}}
\def\Picardfunctor{\mathrm{Pic}}
\def\Deformationcategory{\mathcal{D}\!\mathit{ef}}
\newenvironment{reference}{}{}
\newenvironment{slogan}{}{}
\newenvironment{history}{}{}
\usepackage{cleveref}
\crefname{originallemma}{original lemma}{original lemmas}
\Crefname{originallemma}{Original lemma}{Original lemmas}
\makeatletter
\newlabel{algebra-lemma-characterize-smooth-over-field}{{00TT}{source}{}{sourcepointer-1}{}}
\newlabel{algebra-lemma-characterize-smooth-over-field@cref}{{[originallemma][0][]00TT}{[1][0][]source}{}{}{}}
\makeatother
\crefname{originallemma}{original lemma}{original lemmas}
\Crefname{originallemma}{Original lemma}{Original lemmas}
\makeatletter
\newlabel{algebra-lemma-dimension-at-a-point-finite-type-over-field}{{00OT}{source}{}{sourcepointer-2}{}}
\newlabel{algebra-lemma-dimension-at-a-point-finite-type-over-field@cref}{{[originallemma][0][]00OT}{[1][0][]source}{}{}{}}
\makeatother
\crefname{originallemma}{original lemma}{original lemmas}
\Crefname{originallemma}{Original lemma}{Original lemmas}
\makeatletter
\newlabel{algebra-lemma-minimal-prime-reduced-ring}{{00EU}{source}{}{sourcepointer-3}{}}
\newlabel{algebra-lemma-minimal-prime-reduced-ring@cref}{{[originallemma][0][]00EU}{[1][0][]source}{}{}{}}
\makeatother
\crefname{originalsection}{original section}{original sections}
\Crefname{originalsection}{Original section}{Original sections}
\makeatletter
\newlabel{algebra-section-separability}{{030I}{source}{}{sourcepointer-4}{}}
\newlabel{algebra-section-separability@cref}{{[originalsection][0][]030I}{[1][0][]source}{}{}{}}
\makeatother
\crefname{originallemma}{original lemma}{original lemmas}
\Crefname{originallemma}{Original lemma}{Original lemmas}
\makeatletter
\newlabel{morphisms-lemma-dimension-fibre-at-a-point}{{02FX}{source}{}{sourcepointer-5}{}}
\newlabel{morphisms-lemma-dimension-fibre-at-a-point@cref}{{[originallemma][0][]02FX}{[1][0][]source}{}{}{}}
\makeatother
\crefname{originallemma}{original lemma}{original lemmas}
\Crefname{originallemma}{Original lemma}{Original lemmas}
\makeatletter
\newlabel{morphisms-lemma-finite-type-noetherian}{{01T6}{source}{}{sourcepointer-6}{}}
\newlabel{morphisms-lemma-finite-type-noetherian@cref}{{[originallemma][0][]01T6}{[1][0][]source}{}{}{}}
\makeatother
\crefname{originallemma}{original lemma}{original lemmas}
\Crefname{originallemma}{Original lemma}{Original lemmas}
\makeatletter
\newlabel{properties-lemma-Noetherian-topology}{{01OZ}{source}{}{sourcepointer-7}{}}
\newlabel{properties-lemma-Noetherian-topology@cref}{{[originallemma][0][]01OZ}{[1][0][]source}{}{}{}}
\makeatother
\crefname{originallemma}{original lemma}{original lemmas}
\Crefname{originallemma}{Original lemma}{Original lemmas}
\makeatletter
\newlabel{topology-lemma-Noetherian}{{0052}{source}{}{sourcepointer-8}{}}
\newlabel{topology-lemma-Noetherian@cref}{{[originallemma][0][]0052}{[1][0][]source}{}{}{}}
\makeatother
\title{Free differentials, reducedness and a perfect base field force smoothness}
\author{The Stacks Project Authors}
\date{stacks-project:varieties.tex:local-snapshot-20261011; GNU FDL 1.2 or later}
\begin{document}
\maketitle
\noindent\textit{Editorial note.} All delivered mathematical text is the author's own, copied from the pinned Stacks chapter file varieties.tex; no statement or proof step was rewritten, reordered or removed.\par\medskip
\noindent\textit{Editorial note.} The chapter text immediately before the lemma supplies a counterexample showing that reducedness cannot simply be dropped in characteristic p, that is, the example Spec(k[t]/(t\textasciicircum{}p - a)) with free differentials which is not smooth; that remark lies outside the counted span.\par\medskip
\noindent\textit{Editorial note.} The counted spans carry 8 navigation stubs and no bibliography citation.\par\medskip
\noindent\textit{Editorial note.} The preamble is the chapter's retained minimal preamble (sources/stacks-varieties/retained-preamble.tex lines 1-77), a verbatim excerpt of the native preamble.\par\medskip
\noindent\textit{Editorial note.} The label lemma-char-p-differentials-free-smooth is defined only in varieties.tex, so no cross-chapter restatement had to be weighed for this card.\par\medskip
\noindent\textit{Allowed standard prerequisites (editorial).}\begin{itemize}
\item A finitely generated field extension of a perfect field is separably generated. Algebra, Section algebra-section-separability, invoked at proof lines 4579-4581. Supplies that the module of differentials of the function field is free of rank the transcendence degree.
\item The local ring of a reduced scheme at a generic point of a component is a field. Algebra, Lemma algebra-lemma-minimal-prime-reduced-ring, invoked at proof lines 4576-4578. Supplies the identification of the local ring at eta\_i with kappa(eta\_i).
\item Dimension of a fibre at a point as a transcendence degree. Morphisms, Lemma morphisms-lemma-dimension-fibre-at-a-point, invoked at proof lines 4583-4585. Supplies dim\_\{eta\_i\}(X\_i) = rank of the differentials at eta\_i.
\item Dimension at a point and the criterion for smoothness over a field by comparison with the rank of the differentials. Algebra, Lemmas algebra-lemma-dimension-at-a-point-finite-type-over-field and algebra-lemma-characterize-smooth-over-field, invoked at proof lines 4590-4593. Supplies the final identification of dim\_x(X) with rank\_x(Omega) and hence smoothness at x.
\end{itemize}\medskip
\ifdefined\hypertarget\hypertarget{sourcepointer-1}{}\fi
\noindent\textit{Source reference.} Original lemma 00TT: sources/stacks-chow/tags.txt, source lines 1087--1087. Pointer outside the counted span, invoked once for context; the load-bearing argument of the counted statement is the author's proof reproduced in full in the proof span.\par\medskip
\ifdefined\hypertarget\hypertarget{sourcepointer-2}{}\fi
\noindent\textit{Source reference.} Original lemma 00OT: sources/stacks-chow/tags.txt, source lines 907--907. Pointer outside the counted span, invoked once for context; the load-bearing argument of the counted statement is the author's proof reproduced in full in the proof span.\par\medskip
\ifdefined\hypertarget\hypertarget{sourcepointer-3}{}\fi
\noindent\textit{Source reference.} Original lemma 00EU: sources/stacks-chow/tags.txt, source lines 548--548. Pointer outside the counted span, invoked once for context; the load-bearing argument of the counted statement is the author's proof reproduced in full in the proof span.\par\medskip
\ifdefined\hypertarget\hypertarget{sourcepointer-4}{}\fi
\noindent\textit{Source reference.} Original section 030I: sources/stacks-chow/tags.txt, source lines 3949--3949. Pointer outside the counted span, invoked once for context; the load-bearing argument of the counted statement is the author's proof reproduced in full in the proof span.\par\medskip
\ifdefined\hypertarget\hypertarget{sourcepointer-5}{}\fi
\noindent\textit{Source reference.} Original lemma 02FX: sources/stacks-chow/tags.txt, source lines 3208--3208. Pointer outside the counted span, invoked once for context; the load-bearing argument of the counted statement is the author's proof reproduced in full in the proof span.\par\medskip
\ifdefined\hypertarget\hypertarget{sourcepointer-6}{}\fi
\noindent\textit{Source reference.} Original lemma 01T6: sources/stacks-chow/tags.txt, source lines 2363--2363. Pointer outside the counted span, invoked once for context; the load-bearing argument of the counted statement is the author's proof reproduced in full in the proof span.\par\medskip
\ifdefined\hypertarget\hypertarget{sourcepointer-7}{}\fi
\noindent\textit{Source reference.} Original lemma 01OZ: sources/stacks-chow/tags.txt, source lines 2212--2212. Pointer outside the counted span, invoked once for context; the load-bearing argument of the counted statement is the author's proof reproduced in full in the proof span.\par\medskip
\ifdefined\hypertarget\hypertarget{sourcepointer-8}{}\fi
\noindent\textit{Source reference.} Original lemma 0052: sources/stacks-chow/tags.txt, source lines 194--194. Pointer outside the counted span, invoked once for context; the load-bearing argument of the counted statement is the author's proof reproduced in full in the proof span.\par\medskip
\begin{lemma}
\label{lemma-char-p-differentials-free-smooth}
Let $k$ be a field. Let $X$ be a scheme over $k$.
Assume
\begin{enumerate}
\item $X$ is locally of finite type over $k$,
\item $\Omega_{X/k}$ is locally free,
\item $X$ is reduced, and
\item $k$ is perfect.
\end{enumerate}
Then the structure morphism $X \to \Spec(k)$ is smooth.
\end{lemma}

\begin{proof}
Let $x \in X$ be a point. As $X$ is locally Noetherian (see
Morphisms, Lemma \ref{morphisms-lemma-finite-type-noetherian})
there are finitely many irreducible components
$X_1, \ldots, X_n$ passing through $x$ (see
Properties, Lemma \ref{properties-lemma-Noetherian-topology} and
Topology, Lemma \ref{topology-lemma-Noetherian}).
Let $\eta_i \in X_i$ be the generic point. As $X$ is reduced we have
$\mathcal{O}_{X, \eta_i} = \kappa(\eta_i)$, see
Algebra, Lemma \ref{algebra-lemma-minimal-prime-reduced-ring}.
Moreover, $\kappa(\eta_i)$ is a finitely generated field extension
of the perfect field $k$ hence separably generated over $k$ (see
Algebra, Section \ref{algebra-section-separability}).
It follows that $\Omega_{X/k, \eta_i} = \Omega_{\kappa(\eta_i)/k}$
is free of rank the transcendence degree of $\kappa(\eta_i)$ over $k$. By
Morphisms, Lemma \ref{morphisms-lemma-dimension-fibre-at-a-point}
we conclude that $\dim_{\eta_i}(X_i) = \text{rank}_{\eta_i}(\Omega_{X/k})$.
Since $x \in X_1 \cap \ldots \cap X_n$ we see that
$$
\text{rank}_x(\Omega_{X/k}) = \text{rank}_{\eta_i}(\Omega_{X/k}) = \dim(X_i).
$$
Therefore $\dim_x(X) = \text{rank}_x(\Omega_{X/k})$, see
Algebra, Lemma \ref{algebra-lemma-dimension-at-a-point-finite-type-over-field}.
It follows that $X \to \Spec(k)$ is smooth at $x$ for example by
Algebra, Lemma \ref{algebra-lemma-characterize-smooth-over-field}.
\end{proof}

\end{document}
